\VerdierPhysicalPageAnchor{0023}%
$\mathcal A$ has enough injective objects. The total derived functor therefore
makes it possible, in particular, to reconstruct the ``right hyperderived
functors of $F$''; but it provides more information. For example, when one has
an additive functor $G:\mathcal B\to\mathcal C$ between abelian categories,
one can compose the total derived functors:

\VerdierSemanticAnchor{intro-3-derived-composition}
\[
\mathrm R G\mathbin{\cdot}\mathrm R F:
\mathbf D^*(\mathcal A)\longrightarrow\mathbf D(\mathcal C).
\]
In the usual cases, the functor
$\mathrm R G\mathbin{\cdot}\mathrm R F$ is canonically isomorphic to the
total derived functor of $GF$:

\VerdierSemanticAnchor{intro-3-derived-composition-isomorphism}
\[
\mathrm R(GF)\xrightarrow{\sim}
\mathrm R G\mathbin{\cdot}\mathrm R F,
\]
a composition formula that expresses in our language
\emph{the spectral sequence of composite functors}, and is moreover more
precise, since it is not confined to giving certain information about the
cohomology objects of $\mathrm R GF(X)$
($X$ an object of $\mathbf D^*(\mathcal A)$), but completely describes the
complex ``up to quasi-isomorphism'' that gives rise to these cohomology
objects; knowing this complex makes it possible, for example, to compose once
again with another total derived functor.

\VerdierSemanticAnchor{intro-4}
\noindent\textbf{4.}\quad
Generally speaking, in this formalism one is led to work directly with
complexes ``up to quasi-isomorphism'', which one takes care not to replace by
their cohomology objects alone. (Although most often the ``invariants'' sought
at the end of a calculation or an argument are the cohomology objects of the
complexes, this is not always so: the differentials of spectral sequences are
also interesting ``invariants''.)

\VerdierSemanticAnchor{intro-4-derived-relations}
At the level of complexes up to quasi-isomorphism, or at the level of total
derived functors, or more generally at the level of functors between derived
categories, one establishes various relations (isomorphism relations,
adjunction relations for duality-type theorems, etc.). By purely automatic
procedures (spectral sequences), these relations will provide information
about the cohomology objects of the complexes, about the natural filtrations
with which those objects are equipped, etc.

\VerdierSemanticAnchor{intro-4-technical-intermediaries}
Even when one is interested only in the values of the traditional derived
functors, \emph{i.e.} only in the cohomology objects, relations in derived
categories are often indispensable technical intermediaries for studying
these functors and formulating certain important relations for them.

\VerdierSemanticAnchor{intro-4-simpler-relations}
Finally, an essential advantage is that the relations established in derived
categories are in appearance and in fact simpler than the relations that one
